% Generated draft: requires mathematical and source review.
\documentclass[11pt]{article}
\usepackage[margin=25mm]{geometry}
\usepackage{fontspec}
\setmainfont{Latin Modern Roman}
\usepackage{amsmath,amssymb,mathtools,unicode-math}
\setmathfont{Latin Modern Math}
\usepackage{enumitem,needspace,xurl,hyperref}
\setlength{\parindent}{0pt}
\setlength{\parskip}{5pt}
\setlength{\emergencystretch}{3em}
\newcommand{\N}{\mathbb{N}}
\newcommand{\Z}{\mathbb{Z}}
\newcommand{\Q}{\mathbb{Q}}
\newcommand{\R}{\mathbb{R}}
\newcommand{\C}{\mathbb{C}}
\newcommand{\F}{\mathbb{F}}
\newcommand{\A}{\mathbb{A}}
\newcommand{\PP}{\mathbb{P}}
\newcommand{\E}{\mathbb{E}}
\newcommand{\eps}{\varepsilon}
\newcommand{\dd}{\,\mathrm{d}}
\begin{document}
\section*{UnsolvedMath ID: 11000037}
\textbf{Source problem code:} catalogue entry
\subsection{Definitions and mathematical statement}
Let \(\Sigma_g\) be a closed oriented genus-\(g\) surface and let
\[\mathcal I_g=\ker\bigl(\operatorname{Mod}_g\to\operatorname{Sp}(H_1(\Sigma_g;\Z))\bigr)\]
be its Torelli group. The stated conjecture is
\[d(\mathcal I_g)=0.\]
The source then asks, as a stronger goal, to determine \(d_{\log}(\mathcal I_g)\). Neither \(d\) nor \(d_{\log}\) is defined in this excerpt, so their values are not interpreted as dimensions, stretch factors or any particular density without the missing definition. The genus quantifier for this displayed conjecture also depends on preceding context.

The source says the conjecture would imply \(d(\mathcal I_g(m))=0\) for every \(m\geq2\), and reports exponential growth of \(\mathcal I_g(m)\) for every \(g\geq2,m\geq1\). It asks to compare these exponential growth rates and to quantify how commonly mapping classes act trivially at increasing depths of the lower central series of \(\pi_1\Sigma_g\).

To identify the filtration issue explicitly, put \(\pi=\pi_1\Sigma_g\), \(\gamma_1\pi=\pi\), and \(\gamma_{r+1}\pi=[\pi,\gamma_r\pi]\). A common Johnson-filtration convention is
\[\ker\bigl(\operatorname{Mod}_g\to\operatorname{Out}(\pi/\gamma_{m+1}\pi)\bigr),\]
whose first term is Torelli. This is a candidate interpretation of \(\mathcal I_g(m)\), not a confirmation of the source's indexing or pointed-versus-outer action convention. Triviality on a nilpotent quotient is not silently equated with pointwise triviality on the lower-central subgroups themselves.

Likewise an ambient word-growth comparison, after choosing a finite generating set \(S\) of \(\operatorname{Mod}_g\), could use \(|\mathcal I_g(m)\cap B_S(R)|\), its exponential rate, and its ratio to \(|B_S(R)|\). These illustrate data that a frequency question must specify, not definitions of the missing \(d,d_{\log}\). Intrinsic finite-generator growth is a separate convention and is not automatically available for every filtration subgroup. The source does not state the sampling metric, normalization or order of limits in depth and radius.
\subsection{Short English statement}
The conjecture says the source's invariant \(d\) vanishes on the Torelli group, and asks for its logarithmic refinement and comparisons across the lower-central/Johnson filtration. The density and growth conventions are missing and cannot be replaced by dilatation questions.
\subsection{Sources}
\begin{itemize}
\item Local UnsolvedMath registry, permanent ID \(11000037\), exact source code \texttt{catalogue entry}.
\item Farb, \emph{Problems on Mapping Class Groups and Related Topics}, Conjecture 3.16, PDF page 36: \url{https://www.math.uchicago.edu/~farb/papers/mcgbook.pdf}. Registry access 2026-07-29; pdf-font-structure extraction.
\item The research summary mentions Torelli/dilatation literature without an exact citation and expressly says the definition of \(d\) was unavailable.
\end{itemize}
\end{document}
